\documentclass[11pt]{article}
\usepackage[utf8]{inputenc}
%\usepackage[T1]{fontenc}
\usepackage{amssymb}
\usepackage{amsmath}
\usepackage{enumerate}
\usepackage{fullpage}
\usepackage{polski}  
\usepackage{indentfirst} 
\usepackage[pdftex]{graphicx}
\usepackage{multirow}
\usepackage{placeins}
\usepackage{multirow}

\begin{document}

\begin{tabular}{|p{2.8cm}|cc|l|}
\hline
	Wydział:EAIE & Monika Zygmunt & Maciej Janicki & Rok I (2011/2012) \\
        \cline{4-4}
        Kierunek:EiT & Jan Kasperski  & Łukasz Dubiel & Grupa 6\\
        \cline{4-4}
        & & (bambucha14@gmail.com) & Zespół D\\
\hline
	Data wykonania 2 kwietnia 2012&
        \multicolumn{3}{|p{12cm}|}{Laboratorium Metrologi: Cw 1 i 2
	Wprowadzenie do obsługi przyrządów pomiarowych 
	oraz analiza błędów i niepewności pomiarowych} \\
        
\hline
	Zaliczenie & Podpis prowadzącego & & Uwagi\\
\hline
\end{tabular}
\section{Obserwacja i badanie działanie toru przetwarzania A/C-C/A}
\subsection{Schemat}

\begin{center}
\includegraphics[scale=0.3]{1.pdf}
\end{center}

\subsection{Wyniki pomiarów}
\begin{description}
\item[Amplituda sygnału wejściowego] $U_{mAC} = 5 V $ 
\item[Składowa stała sygnału wejściowego] $U_{0AC} = 0 V $
\item[Amplituda sygnału wyjściowego] $U_{mCA} = 4,32 V $
\item[Ilość stanów przyjmowanych przez tor przetwarzania] $N = 16 $
\item[Rozdzielczość przetwornika] $ N_b = 4 $ 
\item[Częstotliwosć graniczna toru przetwarzania] $ f_{max} = 50kHz $
\end{description}

\begin{center}
\begin{tabular}{|c|c|c|}
\hline
$U_{mAC} [V]$ & $U_{mCA} [V] $ & $N$ \\
\hline
 2 & 2 & 4 \\
\hline
 4 & 3,30 & 6 \\
\hline
 6 & 5,8 & 10 \\
\hline
 8 & 7,84 & 13 \\
\hline
 10 & 9,68 & 16 \\
\hline
 12 & 9,68 & 16 \\
\hline
\end{tabular}
\end{center}

\subsection{Wnioski}

Tor wpływa na propagację sygnału. Częściowo działa jak filtr dolno przepustowy. Można określić tłumienie tego czwórnika. Jest mało odporny na szumy. Problem również stanowi znaczna ilość elementów potrzenych do zwiększenia rozdzielczości przetwornika A/C (rośnie wykładniczo). 

Na częstotliwość graniczną wpływa bardziej przetwornik C/A która jest mniejsza niż częstotliwość maksymalna.

\subsection{Wykaz przyrządów i użytych elementów}
\begin{itemize}
\item{Generator Rigol DG1022}
\item{Oscyloskop Rigol DS1052E}
\item{Tor przetwarzania AC/CA}
\end{itemize} 

\newpage
\section{Obserwacja przebiegów w wybranych punktach przetwornika A/C z dwukrotnym całkowaniem}

\subsection{Schemat}

\begin{center}
\includegraphics[scale=0.3]{4.pdf}
\end{center}

\subsection{Przebiegi napięcia całkowania}

\begin{center}
\includegraphics[scale=0.7]{pod.pdf}
\end{center}

\subsection{Wyniki pomiarów}
\begin{tabular}{|c|c|c|c|}
\hline
Napięcie wejściowe & Czas pierszego całkowania & Czas drugiego całkowania& Maksymalne napięcie \\
$U_{I}[V]$ & $T_1[ms]$ & $T_2[ms]$ & $U_{max}[V]$ \\
\hline
0,9 & 20 & 10 & 15,4 \\
\hline
2,1 & 20 & 20 & 15,6 \\
\hline
0,1 & 20 & 5 & 15,6 \\
\hline
\end{tabular}

\section{Wnioski}
Przetworniki z podójnym całkowaniem są odporne na zakłucenia o określonej cześtotliwości
(oraz kolejne ich harmonicznych).
Pozowala to na usunięcie zakłuceń z sicie energetycznej.
Problematyczny staje się długi czas marwy 
(zdecydowanie większy niż w przypadku przetwonika równoległego bezpośredniego) jednak uzyskuje się znacznie większą dokładność. 

Im mniejsza amplituda sygnału tym sygnał może być próbowany z większą częstotliwością.

\section{Wykaz przyrządów}
\begin{itemize}
\item{Generator Rigol DG1022}
\item{Oscyloskop Rigol DS1052E}
\item{Całkowy przetwornik }
\end{itemize} 


\end{document}
